The numerical experiments presented in this chapter provide a
controlled environment for investigating the proposed cost functional
and optimisation method. The terrain model is intentionally simplified
and should therefore be viewed as a numerical sandbox rather than a
realistic railway or road planning model.

\section{Limitations}

The discretisation provides only an approximation to the continuous
problem, with the curvature term in particular represented using a
finite-difference substitute. Simulated annealing also provides no
guarantee of finding a global minimiser, and the results may depend on
the initial path, random choices and algorithm parameters.

The weighting parameters $(w_c,\lambda,\mu)$ and the annealing
parameters were selected largely through experimentation. The
cost-weighted segmented shift move also contains heuristic choices,
meaning that different parameter values or neighbourhood constructions
may produce different results. A more systematic investigation would
therefore be required to establish the robustness of the numerical
results.

Finally, the terrain considered is relatively simple. More complicated
terrain, obstacles and higher spatial resolutions would increase the
computational cost and provide a more demanding test of the method.

\section{Further Work}

A natural extension would be a systematic investigation of the
parameter space to better understand the trade-off between terrain
cost, curvature and path length. The numerical implementation could
also be refactored using the lessons learned during its development,
reducing the reliance on problem-specific heuristics.

The terrain model could be extended to include multiple terrain
features and impassable regions, where the corresponding areas are
removed from the admissible domain $\Omega$. Introducing elevation
would further extend the problem to three dimensions,
\[
\gamma : [0,1] \rightarrow \mathbb{R}^3,
\]
allowing gradients, ramps and tunnelling to be considered.

A further long-term objective would be to apply the method to
real-world geographical data. Mapping and terrain data could be
converted into a discretised cost field, allowing the optimisation
method to be tested on realistic route-planning problems.

\section{Concluding Remarks}

Despite these limitations, the numerical experiments demonstrate that
including terrain, curvature and length costs can produce non-trivial
paths for which the straight-line route is not necessarily optimal.
The experiments therefore provide numerical evidence supporting the
motivation for the model developed in this dissertation. They also
illustrate the complementary roles of the direct method, which
provides a theoretical framework for the existence of minimisers, and
numerical optimisation, which allows their approximate form to be
investigated computationally.

